\documentclass[12pt,a4paper]{article}
\usepackage[utf8]{inputenc}
\usepackage[T2A]{fontenc}
\usepackage[bulgarian]{babel}
\usepackage{graphicx}
\usepackage{geometry}
\usepackage{enumitem}
\usepackage{xcolor}
\usepackage{amsmath,amssymb}
\usepackage{tikz}
\usetikzlibrary{arrows.meta,calc,positioning}
\usepackage[most]{tcolorbox}
\tcbuselibrary{skins,breakable}

\definecolor{tocgreen}{RGB}{0,110,40}
\definecolor{defblue}{RGB}{20,70,140}
\definecolor{thmorange}{RGB}{190,90,10}
\definecolor{exgray}{RGB}{90,90,90}
\definecolor{probgreen}{RGB}{0,110,40}
\definecolor{defgreen}{RGB}{20, 110, 75}

\definecolor{greenaccent}{HTML}{2E7D32}
\definecolor{blueaccent}{HTML}{000080}

\usepackage[colorlinks=true,linkcolor=tocgreen,urlcolor=tocgreen,citecolor=tocgreen]{hyperref}
\geometry{margin=2cm}

\newcommand{\vv}[1]{\overrightarrow{#1}}

\newtcolorbox{defbox}[1][Определение]{
  colback=defgreen!5, colframe=defgreen, coltitle=white,
  fonttitle=\bfseries, title=#1, boxrule=0.8pt, arc=2pt,
  left=6pt, right=6pt, top=4pt, bottom=4pt, breakable
}
\newtcolorbox{thmbox}[1][Теорема]{
  colback=thmorange!6, colframe=thmorange, coltitle=white,
  fonttitle=\bfseries, title=#1, boxrule=0.8pt, arc=2pt,
  left=6pt, right=6pt, top=4pt, bottom=4pt, breakable
}
\newtcolorbox{exbox}[1][Пример]{
  colback=exgray!5, colframe=exgray, coltitle=white,
  fonttitle=\bfseries\itshape, title=#1, boxrule=0.6pt, arc=2pt,
  left=6pt, right=6pt, top=4pt, bottom=4pt, breakable
}
\newtcolorbox{probbox}[1][Задача]{
  colback=probgreen!4, colframe=probgreen, coltitle=white,
  fonttitle=\bfseries, title=#1, boxrule=0.7pt, arc=2pt,
  left=6pt, right=6pt, top=4pt, bottom=4pt, breakable
}

\newcommand{\seclabel}{lesson11}

\newcommand{\qtag}[1]{\texttt{\color{greenaccent}\small#1}}
\newcommand{\qtagb}[1]{\texttt{\color{blueaccent}\small#1}}

\begin{document}

\begin{titlepage}
\centering
\vspace*{2cm}
{\large МИНИСТЕРСТВО НА ОБРАЗОВАНИЕТО И НАУКАТА\par}
\vspace{-0.2cm}
\noindent\rule{0.8\textwidth}{0.5pt}
\vspace{0.3cm}
{\large ПОДГОТОВКА ЗА ДЪРЖАВЕН ЗРЕЛОСТЕН ИЗПИТ\par}
\vspace{-0.4cm}
\noindent\rule{0.8\textwidth}{0.5pt}
\vspace{0.8cm}
\vspace{1.5cm}
{\Huge\bfseries МАТЕМАТИКА\par}
\vspace{0.6cm}
{\Large Подготовка за ДЗИ -- 12. клас\par}
\vspace{0.6cm}
{\Large Сборник с тестове и примерни изпити\par}
\vspace{2.4cm}
\noindent\rule{1.0\textwidth}{0.5pt}
\vspace{0.4cm}
{\Large\bfseries Учебна година\par}
{\Large\bfseries 2026/2027 \par}
\vspace{-0.2cm}
\noindent\rule{1.0\textwidth}{0.5pt}
{\Large\bfseries Автор:\par}
{\large Катерина Димитрова\par}
\noindent\rule{1.0\textwidth}{0.5pt}
\vfill
{\large $\bullet$ София $\bullet$ \par}
{\large 2026 \par}
\end{titlepage}

\newpage
\begin{center}
{\LARGE\bfseries\color{tocgreen} Съдържание}
\end{center}
\vspace{0.5cm}

{\color{tocgreen}
\noindent{\large\bfseries Модул I -- Геометрия}\par
\vspace{0.2cm}
\noindent{\bfseries I. Вектори и координати}\par
\begin{enumerate}[label=1.\arabic*.,leftmargin=1.5cm]
\item \hyperref[\seclabel]{Линейна зависимост и независимост на вектори в равнината и пространството} \dotfill \hyperref[\seclabel]{5}
\item Векторна база в равнината и пространството \dotfill 9
\item Скаларно произведение на два вектора. Приложение \dotfill 11
\item Координати на вектор в равнинна правоъгълна координатна система \dotfill 14
\item Операции с вектори, зададени с координати \dotfill 16
\end{enumerate}
\noindent Вектори и координати -- Тест 1 и Тест 2 \dotfill 19

\vspace{0.3cm}
\noindent{\bfseries II. Аналитична геометрия в равнината}\par
\begin{enumerate}[label=2.\arabic*.,leftmargin=1.5cm]
\item Уравнение на права \dotfill 21
\item Взаимно положение на две прави \dotfill 26
\item Приложение на векторите в аналитичната геометрия за решаване на триъгълник \dotfill 30
\item Нормално уравнение на окръжност \dotfill 33
\item Канонично уравнение на елипса, хипербола и парабола \dotfill 35
\end{enumerate}
\noindent Аналитична геометрия в равнината -- Тест 1 и Тест 2 \dotfill 38

\vspace{0.3cm}
\noindent{\bfseries III. Стереометрия}\par
\begin{enumerate}[label=3.\arabic*.,leftmargin=1.5cm]
\item Първични понятия и аксиоми в стереометрията. Успоредност в пространството \dotfill 40
\item Перпендикулярност в пространството \dotfill 44
\item Перпендикуляр и наклонена \dotfill 48
\item Двустенен ъгъл. Перпендикулярност на две равнини \dotfill 50
\item Многостен \dotfill 52
\item Сечение на многостен с равнина \dotfill 55
\item Построяване на сечение с равнина \dotfill 59
\item Ос на кръстосани прави \dotfill 63
\item Ротационни тела \dotfill 68
\end{enumerate}
\noindent Стереометрия -- Тест 1 и Тест 2 \dotfill 74\par
\noindent Годишен преговор \dotfill 83\par
\noindent Проекти \dotfill 86

\vspace{0.5cm}
\noindent{\large\bfseries Модул II -- Елементи на математическия анализ}\par
\vspace{0.2cm}
\noindent{\bfseries I. Полиноми на една променлива}\par
\begin{enumerate}[label=1.\arabic*.,leftmargin=1.5cm]
\item Определение. Операции с полиноми \dotfill 90
\item Теорема на Безу. Схема на Хорнер \dotfill 93
\item Нули на полиноми \dotfill 97
\item Рационални корени на уравнение с цели коефициенти \dotfill 99
\item Решаване на уравнения и неравенства от по-висока степен \dotfill 102
\end{enumerate}
\noindent Полиноми на една променлива -- Тест 1 и Тест 2 \dotfill 106

\vspace{0.3cm}
\noindent{\bfseries II. Числови редици}\par
\begin{enumerate}[label=2.\arabic*.,leftmargin=1.5cm]
\item Метод на математическата индукция \dotfill 107
\item Нютонов бином \dotfill 109
\item Числови редици \dotfill 112
\item Теореми за граници на редици \dotfill 115
\item Сума на безкрайно намаляваща геометрична прогресия \dotfill 123
\end{enumerate}
\noindent Числови редици -- Тест 1 и Тест 2 \dotfill 125

\vspace{0.3cm}
\noindent{\bfseries III. Функции. Непрекъснатост и диференцируемост}\par
\begin{enumerate}[label=3.\arabic*.,leftmargin=1.5cm]
\item Функция. Начини на задаване \dotfill 127
\item Съставна функция \dotfill 129
\item Граница на функция \dotfill 130
\item Теореми за граница на функция \dotfill 132
\item Основни граници \dotfill 136
\item Непрекъснатост \dotfill 139
\item Теореми за непрекъснатост \dotfill 142
\item Производна на функция \dotfill 144
\item Връзка между непрекъснатост и диференцируемост \dotfill 150
\end{enumerate}
\noindent Функции. Непрекъснатост и диференцируемост -- Тест 1 и Тест 2 \dotfill 152\par
\noindent Проекти \dotfill 154\par
\noindent Годишен преговор \dotfill 159

\vspace{0.3cm}
\noindent Отговори -- Модул I \dotfill 163\par
\noindent Отговори -- Модул II \dotfill 169
}

\vspace{0.5cm}

\newpage

% УРОК 1.1
\phantomsection\label{\seclabel}

\begin{tcolorbox}[
    blanker,
    top=10pt, bottom=10pt,
    borderline architectural={1.5pt}{0pt}{tocgreen}
]
{\Large\bfseries\color{tocgreen!90!black} 1.1. Линейна зависимост и независимост на вектори\\ в равнината и в пространството}
\end{tcolorbox}
\vspace{-0.7cm}

\subsection*{А. Теория}

\subsubsection*{1) Дефиниция за вектор}

\begin{defbox}[Дефиниция 1]
Отсечка, при която единият край е избран за начало, а другият -- за край, се нарича \emph{насочена отсечка} или \emph{вектор}. \\ Записваме \qtag{$\vv{AB}$}, като $A$ е началото, а $B$ -- краят на насочената отсечка. \\ Вектор може да се бележи и с малка латинска буква: \qtag{$\vv{AB}=\vec{a}$}.
Въвеждаме следните понятия:
\begin{itemize}
    \vspace{-0.2cm}
    \item[\textbf{--}] \textbf{Нулев вектор:} Ако началото и краят на вектора съвпадат \qtag{($B=A$)}, той се нарича нулев вектор и се бележи с \qtag{$\vv{AA}=\vec{0}$}.
    \vspace{-0.2cm}
    \item[\textbf{--}] \textbf{Дължина на вектор:} Дължината на отсечката $AB$ се нарича дължина (или модул) на вектора $\vv{AB}$ и се бележи с \qtag{$|\vv{AB}|$} или \qtag{$|\vec{a}|$}.
\end{itemize}
\end{defbox}

\begin{center}
\begin{tikzpicture}[>=Latex, thick, scale=1.2]
    \draw[step=0.5cm, gray!15, very thin] (-2.3,-0.7) grid (5.5,2.5);
    \node[anchor=north west, fill=white, inner sep=2pt, text=black!60, font=\small\rmfamily\bfseries] at (-2.0,1.9) {Пример:};
    \draw[->, gray!60, thin] (-0.5,0) -- (5,0) node[right, scale=0.8, text=black!70] {$x$};
    \draw[->, gray!60, thin] (0,-0.3) -- (0,1.7) node[above, scale=0.8, text=black!70] {$y$};
    \coordinate (A) at (0,0);
    \coordinate (B) at (4,1.2);
    \draw[dashed, gray!40, thin] (4,0) -- (B);
    \draw[dashed, gray!40, thin] (0,1.2) -- (B);
    \draw[-, defgreen!15, line width=4pt, line cap=round] (A) -- (B);
    \draw[->, defgreen, line width=1.6pt, line cap=round] (A) -- (B);
    \filldraw[black] (A) circle (2pt);
    \filldraw[white] (A) circle (0.8pt);
    \node[below left=3pt, black!80] at (A) {$A$};
    \node[above right=2pt, black!80] at (B) {$B$};
    \path (A) -- (B) node[midway, above left=2pt, color=defgreen!80!black, font=\boldmath] {$\vec{a}$};
\end{tikzpicture}
\end{center}

\end{document}
